\subsection{非共鳴主chartの多項式逆行列と四つの残余曲面}
\begin{proposition}
非退化正top・対称secant負top・正Pfaffianの $q=1$ 主chartで
\begin{align*}
 L={}&\rho^5+\rho^4+2\rho^3s^2+2\rho^3t^2-2\rho^2s^2
       +2\rho^2t^2+2\rho t^2-\rho+2t^2-1
\end{align*}
と置く。
\[
 s t(\rho^2-2t-1)(\rho^2+2t-1)L\ne0
\]
ならば正常三次核は零であり、H3対象はない。
従って既に閉じた境界・共鳴面を除く未知は
\[
 s=0,\qquad t=0,\qquad\rho^2-2t-1=0,\qquad L=0
\]
の四曲面上に限られる。各曲面は $a\ne0$、$\rho^2+2t-1\ne0$ に制限して調べればよい。
\end{proposition}
\begin{proof}
前節の八つの多項式四次汎関数から、全正自由三次32語に対する正常行列 $N$ を作る。
その各行は元の全三次関係を消す。三次関係との一次式の左右積が四次関係に属し、
四次汎関数がそれらを消すためである。
従って $N$ は実際の正三次商を通る正常条件である。

入力
\[
 \texttt{cubic\_order2\_22\_secant\_main\_open\_input.json}
\]
は $N$ の十行十列部分行列 $C$ と、多項式係数の十次正方行列 $X$ を指定する。
係数環 $\mathbb Q[\rho,s,t]$ 上で
\[
 C X=dI_{10},
\]
\[
 d=\rho^3 s t^2(\rho-1)^5(\rho+1)^6
       (\rho^2-2t-1)(\rho^2+2t-1)L
\]
が恒等的に成立する。
従って主張した開集合で $C$ は可逆、$N$ のrankは10以上となる。
一方、AS4候補の正三次商次元は10であるから、正常核は零でなければならない。
この証明はgeneric三次基底を特殊値で使わず、全32自由語上のrankと元関係恒等式だけを用いる。
また大きな多変数行列式の再計算も不要である。

\texttt{verify\_cubic\_order2\_22\_secant\_main\_open.py}
は元288四次関係を消す汎関数恒等式と $CX=dI_{10}$、$d$ の因子表示を直接再生する。
この必要条件を、四残余曲面の排除や必要atlasのAS4十分性と取り違えてはならない。
\end{proof}
